\section*{Chapitre \ref{ch:m2:credit-default-swaps} --- Credit default swaps}

\begin{solution}{exo:m2:credit-default-swaps:1}
25\,000~USD par trimestre ($10 \text{ millions} \times 1\% / 4$)~; après le défaut, $10
\text{ millions} \times (100 - 40)\% = \text{USD}~6$ millions, et les primes cessent.
\end{solution}

\begin{solution}{exo:m2:credit-default-swaps:2}
$\lambda \approx 0.03 / 0.6 = 5\%$ par an~; $Q(5) = e^{-0.25} = 77.9\%$. C'est une
probabilité neutre au risque, et non une prévision.
\end{solution}

\begin{solution}{exo:m2:credit-default-swaps:3}
Savoir si un \omterm{def:m2:credit-default-swaps:event}{événement de crédit} s'est produit, et quand~; savoir s'il faut organiser une enchère~; quelles
obligations sont livrables~; savoir si l'entité a un successeur, et lequel. Il répond aussi
aux questions d'interprétation du contrat.
\end{solution}

\begin{solution}{exo:m2:credit-default-swaps:4}
Le \omterm{def:m2:credit-default-swaps:std}{paiement initial} (upfront) est de $-2.21\%$~: l'\omterm{def:m2:credit-default-swaps:cds}{acheteur de protection} reçoit environ 221\,000~USD, car il
va payer un coupon de 100 points de base pour une protection valant 50.
\end{solution}

\begin{solution}{exo:m2:credit-default-swaps:5}
La compensation exige des contrats identiques~: deux transactions sur la même \omterm{def:m2:credit-default-swaps:cds}{entité de référence} et la même échéance mais avec
des coupons différents laissent un flux résiduel de primes qui ne peut être annulé. Avec un
\omterm{def:m2:credit-default-swaps:std}{coupon standard}, la différence est payée une fois, à l'avance, et les jambes courantes
sont identiques, si bien que des transactions opposées s'annulent et qu'une contrepartie centrale
peut détenir une position nette par entité et par échéance.
\end{solution}

\begin{solution}{exo:m2:credit-default-swaps:6}
Il livre l'obligation à 30 et reçoit le pair~: 10~millions de dollars pour des obligations valant
3~millions de dollars, soit un paiement de 7~millions de dollars. Le \omterm{def:m2:credit-default-swaps:final}{prix final d'enchère} doit refléter le
\omterm{def:m2:bond-futures:ctd}{moins cher à livrer}, soit environ 30, puisque les vendeurs à l'enchère le livrent.
\end{solution}

\begin{solution}{exo:m2:credit-default-swaps:7}
1.6667~\%, 5.0001~\% et 13.335~\%, contre 1.6667~\%, 5.0000~\% et 13.333~\% issus du
triangle. Le petit écart, qui croît avec le spread, vient du fait que les primes et les
défauts sont versés chaque trimestre plutôt qu'en continu~; en comptant une demi-période de
primes courues, les jambes discrètes coïncident presque avec les continues.
\end{solution}

\begin{solution}{exo:m2:credit-default-swaps:8}
Le spread vaut à peu près l'\omterm{def:m2:credit-default-swaps:hazard}{intensité de défaut} multipliée par la perte en cas de défaut, $\lambda(1-R)$~: à
300 points de base et un \omterm{def:m2:corporate-bonds:covenant}{taux de recouvrement} de 40~\%, l'\omterm{def:m2:credit-default-swaps:hazard}{intensité de défaut} vaut environ 5~\% par an, et non 3~\%. Il
s'agit d'un taux neutre au risque, qui inclut une prime de compensation du risque de défaut et
de l'illiquidité~; les taux de défaut observés sont en général plus faibles.
\end{solution}

\begin{solution}{pb:m2:credit-default-swaps:1}
\textbf{1.} Le meilleur cours acheteur, 10, et le meilleur cours vendeur, 10~: ils se croisent et sont retirés.
\textbf{2.} Sept des treize restants~: cours acheteurs 9.5, 9.5, 9.25, 9.25, 9, 9, 8.875 et
cours vendeurs 10, 10, 10, 10.25, 10.75, 10.875, 11 dont la moyenne vaut 9.804, arrondie à 9.75.
\textbf{3.} Vend pour 5\,692 millions de dollars, achète pour 772 millions de dollars, soit un net à vendre de 4\,920 millions de dollars.
\textbf{4.} Le dealer qui a affiché un cours acheteur de 10, HSBC, au-dessus du point médian, alors que les positions ouvertes étaient
à vendre~: $5\text{ millions} \times (10 - 9.75)\% = \text{USD}~12\,500$.
\textbf{5.} Les \omterm{def:m2:credit-default-swaps:cds}{acheteurs de protection} qui détenaient des obligations demandaient à les vendre au prix final, afin
de finir au pair, et ils pesaient plus que les \omterm{def:m2:credit-default-swaps:cds}{vendeurs de protection} demandant à acheter.
\textbf{6.} Ordres d'achat à cours limité sur les obligations, pour un prix d'au plus $9.75 + 1 = 10.75$.
\textbf{7.} Tailles cumulées de 50, 200, 500, 1\,400 et 2\,800 millions en descendant jusqu'à 8.75, puis
5\,400 millions à 8.625, ce qui couvre 4\,920~: le prix final est de 8.625.
\textbf{8.} 10.75, le cap.
\textbf{9.} 3\,800 millions à 8.625 ne suffisent pas~; avec les 1\,800 millions à 8.5, le
carnet d'ordres atteint 5\,600~: le prix final est 8.5.
\textbf{10.} 1.125 sous le point médian~: les positions ouvertes à vendre étaient importantes, et les cours acheteurs
nécessaires pour l'absorber étaient plus bas.
\textbf{11.} $20 \text{ millions} \times 91.375\% = \text{USD}~18.275$ millions~; ses obligations
valent 1.725 millions de dollars à 8.625.
\textbf{12.} 18.275 millions de dollars au titre du CDS et 1.725 millions de dollars pour les obligations~: 20
millions de dollars, au pair, et plus aucune obligation.
\textbf{13.} 45.6875 millions de dollars.
\textbf{14.} Sur ses 40 millions de dollars nets vendus, 36.55 millions de dollars.
\textbf{15.} 27.4125 millions de dollars sur le CDS et 2.5875 millions de dollars pour les obligations~: 30
millions de dollars au total, pour 30 millions de dollars de valeur nominale, comme s'il avait livré ces obligations.
\textbf{16.} La plus grande partie de la protection était détenue par des dealers qui avaient à la fois acheté et vendu
cette protection~; à l'enchère, chacun ne payait ou ne recevait que sa position nette.
\textbf{17.} Les positions ouvertes à vendre s'élevaient à 4.92~milliards de dollars~; absorber cette offre a exigé
des cours acheteurs inférieurs aux derniers cours du marché obligataire.
\textbf{18.} Les détenteurs de protection sans obligations auraient dû en acheter, faisant monter le
prix et réduisant leur paiement, et les contrats sur la même entité auraient été réglés à
des prix différents.
\textbf{19.} Résultat par acteur~: \emph{l'enchère} fixe le prix final à 8.625~; A reçoit
18.275 millions de dollars et conserve des obligations valant 1.725 (ou 20 millions de dollars sans obligations), B
reçoit 45.6875 millions de dollars, C paie 36.55 millions de dollars, et D paie au total 30 millions de dollars
pour 30 millions de dollars de valeur nominale d'obligations.
\textbf{20.} Elle donne à chaque contrat portant sur une entité un seul prix de recouvrement, fixé à partir de la
demande nette pour ses obligations, de sorte que le règlement-livraison ne dépend plus de la recherche de titres.
\end{solution}

\begin{solution}{iq:m2:credit-default-swaps:1}
Un contrat dans lequel l'\omterm{def:m2:credit-default-swaps:cds}{acheteur de protection} verse une prime courante sur un notionnel au
vendeur, chaque trimestre, jusqu'à l'échéance ou à un \omterm{def:m2:credit-default-swaps:event}{événement de crédit} de l'\omterm{def:m2:credit-default-swaps:cds}{entité de référence}~; après un
\omterm{def:m2:credit-default-swaps:event}{événement de crédit}, le vendeur paie le notionnel multiplié par un moins le prix de recouvrement fixé par
l'enchère. L'acheteur n'a pas besoin de détenir la dette.
\iqlookfor{les deux jambes, le déclencheur, le règlement-livraison.}
\end{solution}

\begin{solution}{iq:m2:credit-default-swaps:2}
Depuis 2009, les contrats portent un coupon fixe (100 ou 500 points de base pour les
entités nord-américaines) et la différence avec le spread de marché est versée en \omterm{def:m2:credit-default-swaps:std}{paiement initial} (upfront). Des jambes courantes
identiques rendent fongibles les transactions sur la même entité et la même échéance~: elles peuvent être compensées,
compressées et admises en chambre de compensation.
\iqlookfor{fongibilité, netting, compensation.}
\end{solution}

\begin{solution}{iq:m2:credit-default-swaps:3}
Avec une \omterm{def:m2:credit-default-swaps:hazard}{intensité de défaut} constante $\lambda$, un \omterm{def:m2:corporate-bonds:covenant}{taux de recouvrement} $R$ et des primes en continu, la jambe de primes est
$s\int e^{-(r+\lambda)t}dt$ et la jambe de protection $(1-R)\lambda\int
e^{-(r+\lambda)t}dt$~; l'égalisation donne $s = \lambda(1-R)$. Des paiements discrets
modifient légèrement ce résultat~; l'\omterm{def:m2:credit-default-swaps:hazard}{intensité de défaut} est neutre au risque.
\iqlookfor{les deux jambes et l'annulation.}
\end{solution}

\begin{solution}{iq:m2:credit-default-swaps:4}
Première étape~: les dealers soumettent des cours acheteur et des cours vendeur pour une taille fixe, ainsi que des demandes de règlement-livraison physique
demandes~; les paires croisées sont retirées et les meilleurs cours sont moyennés pour donner le milieu du
marché intérieur~; les demandes nettes constituent les positions ouvertes. Deuxième étape~: les ordres à cours limité
du côté opposé exécutent les positions ouvertes, par ordre de prix, et le dernier ordre exécuté fixe
le prix final, à un point près du milieu. Elle est nécessaire parce que la protection peut
dépasser les obligations livrables, et les obligations se négocient peu après le défaut~; elle donne un seul prix pour
tous les contrats.
\iqlookfor{les deux étapes, le plafond, la raison.}
\end{solution}

\begin{solution}{iq:m2:credit-default-swaps:5}
La base est le spread du CDS moins le spread de l'obligation par rapport au taux de financement
(pour une obligation à taux variable, son spread~; pour une obligation à taux fixe, un asset swap
ou un Z-spread). Si elle est négative, on achète l'obligation, on la finance en repo et on achète
une protection~: on gagne l'écart et on est couvert contre le défaut. Les risques sont le financement
et les décotes du repo, l'écart entre le prix de l'obligation et le prix d'enchère, les options de
livraison, et la valorisation au prix de marché si la base s'élargit.
\iqlookfor{le montage, le financement et les risques résiduels.}
\end{solution}

\begin{solution}{iq:m2:credit-default-swaps:6}
Conserver les transactions avec leurs conditions standard~; chaque nuit, calibrer par bootstrap les courbes
d'\omterm{def:m2:credit-default-swaps:hazard}{intensité de défaut} par nom et le \omterm{def:m2:corporate-bonds:covenant}{taux de recouvrement} à partir des cotations de marché, valoriser chaque
transaction et agréger le risque par nom et par échéance, en parallèle sur les noms~; en cas
d'\omterm{def:m2:credit-default-swaps:event}{événement de crédit}, suivre les décisions du comité et le résultat de l'enchère, puis calculer le
règlement de chaque transaction au prix final, compenser par contrepartie et par la chambre de
compensation, comptabiliser les paiements~; tester sur les enchères passées.
\iqlookfor{courbes par nom, valorisation en parallèle, processus en cas d'événement.}
\end{solution}
